\chapter{Séries multivariées et cointégration}\label{ch:qm:multivariate-series-and-cointegration}

Les rendements des bons du Trésor américain à deux ans, cinq ans et dix ans ressemblent chacun à des marches aléatoires~: sur cinquante ans de données quotidiennes, aucun des trois ne rejette
une \omterm{def:qm:linear-time-series:unitroot}{racine unitaire} (statistiques de Dickey--Fuller de $-1.30$, $-1.33$ et $-1.31$). Une butterfly de courbe, longue sur le cinq ans et courte sur les ailes, ne
s'éloigne pas de la même façon~: les rendements partagent leurs tendances, et une combinaison des trois est stationnaire. La question que se pose un
desk de relative value est de savoir quelles pondérations la rendent stationnaire, et à quelle vitesse la butterfly qui en résulte revient à la moyenne. La butterfly du manuel pondère les ailes un-deux-un~; les données préfèrent
à mettre 0.41 sur le deux ans et 0.61 sur le dix ans pour chaque unité du cinq ans, et la butterfly ainsi obtenue a une \omterm{def:qm:stochastic-differential-equations:ou}{demi-vie} de 43 jours ouvrés contre 63.
Ce chapitre porte sur plusieurs séries à la fois~: les \omterm{def:qm:multivariate-series-and-cointegration:var}{modèles vectoriels autorégressifs} et l'effet de leurs chocs, la \omterm{def:qm:multivariate-series-and-cointegration:coint}{cointégration} et la forme à
correction d'erreur, les \omterm{def:qm:multivariate-series-and-cointegration:eg}{tests d'Engle--Granger} et de Johansen qui la détectent, et la causalité au sens prédictif.

\section{Modèles vectoriels autorégressifs}

\begin{definition}[Modèle vectoriel autorégressif]\label{def:qm:multivariate-series-and-cointegration:var}
Un \emph{modèle vectoriel autorégressif}\index{modèle vectoriel autorégressif} (vector autoregression) d'ordre $p$, VAR($p$), modélise un vecteur de dimension $k$ comme $Y_t = c + \sum_{i=1}^pA_iY_{t-i} + \varepsilon_t$ avec
$\varepsilon_t$ un \omterm{def:qm:linear-time-series:stationary}{bruit blanc} de covariance $\Sigma$ (Sims, 1980). Chaque équation est estimée par moindres carrés sur les mêmes régresseurs.
\end{definition}

\begin{proposition}[Stationnarité d'un VAR]\label{prop:qm:multivariate-series-and-cointegration:stable}
Le VAR est faiblement stationnaire si et seulement si toute valeur propre de sa matrice compagnon $\begin{psmallmatrix}A_1 & \cdots & A_{p-1} & A_p\\ I & & & 0\\ & \ddots & & \vdots\\
& & I & 0\end{psmallmatrix}$ est de module strictement inférieur à un~; de façon équivalente, $\det(I - \sum_iA_iz^i) \ne 0$ pour $|z| \le 1$.
\end{proposition}

\begin{proof}
On empile $(Y_t, \dots, Y_{t-p+1})$ en un VAR(1) de matrice compagnon $C$. Ses coefficients de moyenne mobile sont $C^h$, qui sont sommables exactement lorsque le
rayon spectral de $C$ est inférieur à un~; les valeurs propres de $C$ sont les inverses des racines du déterminant.
\end{proof}

\begin{definition}[Fonction de réponse impulsionnelle, décomposition de la variance des erreurs de prévision]\label{def:qm:multivariate-series-and-cointegration:irf}
La \emph{fonction de réponse impulsionnelle}\index{fonction de réponse impulsionnelle} donne l'effet sur $Y_{t+h}$ d'un choc égal à un écart type à la date $t$~; lorsque les chocs sont corrélés,
ils sont d'abord orthogonalisés par la factorisation de Cholesky de $\Sigma$, ce qui rend le résultat dépendant de l'ordre des variables. La
\emph{décomposition de la variance des erreurs de prévision}\index{décomposition de la variance des erreurs de prévision} donne la part de la variance de l'erreur de prévision à $h$ pas de chaque variable
attribuable à chacun des chocs orthogonalisés.
\end{definition}

Sur les variations quotidiennes des rendements (en points de base), l'AIC continue de s'améliorer jusqu'aux dix retards autorisés, et le VAR(10) est stable. Les résidus ont des
écarts types de 7.9, 7.6 et 7.1 points de base et sont fortement corrélés (0.90 entre les deux et cinq ans, 0.82 entre les deux et dix ans, 0.94 entre
les cinq et dix ans), donc l'ordre compte. Avec le deux ans placé en premier, un choc d'un écart type sur le deux ans, soit 7.9 points de base, déplace les rendements à deux, cinq
et dix ans de 10.1, 8.3 et 6.7 points de base cumulés sur dix jours (\cref{fig:qm:multivariate-series-and-cointegration:irf}), et il explique
67~\% de la variance de l'erreur de prévision à dix jours du dix ans (le choc du cinq ans pour 21~\%, celui du dix ans lui-même pour 12~\%). Le deux ans mène, dans cet ordre,
par construction~; savoir s'il mène dans le temps est une autre question.

\begin{omfigure}
\begin{tikzpicture}
\begin{axis}[width=10cm,height=5cm,xlabel={jours après le choc},ylabel={réponse cumulée (pb)},
  tick label style={font=\small},label style={font=\small},xmin=0,xmax=10,ymin=0,ymax=11,grid=major,grid style={gray!25},
  legend columns=3,legend style={font=\footnotesize,at={(0.5,-0.3)},anchor=north,draw=none,fill=none,
  /tikz/every even column/.append style={column sep=8pt}},table/col sep=comma]
\addplot[omDef,very thick,mark=*,mark size=1.3pt] table[x=day,y=y2]{figdata/methods/20-multivariate-series-and-cointegration/irf.csv};
\addplot[black,very thick,mark=square*,mark size=1.3pt] table[x=day,y=y5]{figdata/methods/20-multivariate-series-and-cointegration/irf.csv};
\addplot[omThm,very thick,mark=triangle*,mark size=1.6pt] table[x=day,y=y10]{figdata/methods/20-multivariate-series-and-cointegration/irf.csv};
\legend{2 ans,5 ans,10 ans}
\end{axis}
\end{tikzpicture}

\omcaption{Réponses orthogonalisées cumulées des rendements à 2, 5 et 10 ans à un choc d'écart type (7.9 pb) sur le 2 ans, estimé à partir d'un VAR(10) sur
les variations quotidiennes, 1976--2026, ordonnées 2, 5 et 10 ans. Données~: séries FRED DGS2, DGS5 et DGS10 (Board of Governors of the Federal Reserve System,
H.15).}\label{fig:qm:multivariate-series-and-cointegration:irf}
\end{omfigure}

\section{Causalité au sens prédictif}

\begin{definition}[Causalité au sens de Granger]\label{def:qm:multivariate-series-and-cointegration:granger}
Une série $X$ \emph{cause au sens de Granger}\index{cause au sens de Granger} $Y$ si le passé de $X$ améliore la prédiction de $Y$ au-delà du passé de $Y$ et de celui des autres
variables (Granger, 1969)~; dans un VAR, si les coefficients des retards de $X$ dans l'équation de $Y$ ne sont pas tous nuls, testé au moyen d'une statistique $F$.
\end{definition}

Il s'agit d'une antériorité prédictive, non d'une causalité~: un facteur commun qui apparaît plus tôt dans une série que dans l'autre la produit, tout comme un
décalage de datation dans la manière dont les données sont enregistrées. Sur les variations quotidiennes des rendements, chaque couple \omterm{def:qm:multivariate-series-and-cointegration:granger}{cause au sens de Granger} l'autre au seuil de 5~\%, avec $F(10, 12\,532)$ compris entre 1.88 (deux ans vers
dix ans, $p = 0.04$) et 7.16 (deux ans vers cinq ans, $p < 10^{-6}$)~; la causalité du dix ans vers le deux ans donne 2.83 ($p = 0.002$). Avec 12\,574 jours, de minuscules
effets d'avance-retard sont significatifs~; aucun n'est assez grand pour être exploité isolément.

\section{Cointégration et forme à correction d'erreur}

\begin{definition}[Cointégration, vecteur de cointégration, rang de cointégration]\label{def:qm:multivariate-series-and-cointegration:coint}
Les composantes d'un vecteur de processus à \omterm{def:qm:linear-time-series:unitroot}{racine unitaire} sont cointégrées (\emph{cointégration}\index{cointégration}) si une combinaison linéaire $\beta^\top Y_t$ est
stationnaire~; $\beta$ est un \emph{vecteur de cointégration}\index{vecteur de cointégration}. Le nombre de vecteurs de cointégration linéairement indépendants est le
\emph{rang de cointégration}\index{rang de cointégration} $r$, et $k - r$ est le nombre de tendances stochastiques communes.
\end{definition}

\begin{definition}[Modèle vectoriel à correction d'erreur]\label{def:qm:multivariate-series-and-cointegration:vecm}
Un \emph{modèle vectoriel à correction d'erreur}\index{modèle vectoriel à correction d'erreur} (vector error-correction model, VECM) écrit un VAR en différences augmenté d'un terme en niveaux, $\Delta Y_t = \alpha\beta^\top Y_{t-1} +
\sum_{i=1}^{p-1}\Gamma_i\Delta Y_{t-i} + \mu + \varepsilon_t$, avec $\alpha$ et $\beta$ de dimension $k \times r$~: $\beta^\top Y_{t-1}$ est l'écart à l'équilibre
et $\alpha$ la vitesse à laquelle chaque variable le corrige.
\end{definition}

D'après le théorème de représentation de Granger (Engle et Granger, 1987), un système cointégré admet une représentation à correction d'erreur, et réciproquement~: un VAR en niveaux
dont la matrice $\Pi = \sum_iA_i - I$ est de rang réduit $r$ se factorise en $\alpha\beta^\top$. Un VAR en différences seules ignore l'information sur les niveaux~; un VAR en
niveaux la conserve mais la masque. Pour les rendements, $k = 3$~; si leur courbe est portée par des tendances de niveau et de pente avec une courbure stationnaire, $r = 1$ et le
\omterm{def:qm:multivariate-series-and-cointegration:coint}{vecteur de cointégration} est une butterfly de courbe.

\begin{definition}[Test d'Engle--Granger]\label{def:qm:multivariate-series-and-cointegration:eg}
Le \emph{test d'Engle--Granger}\index{test d'Engle--Granger} régresse une variable sur les autres en niveaux et applique un \omterm{def:qm:linear-time-series:df}{test de Dickey--Fuller} aux résidus,
avec des valeurs critiques qui tiennent compte des coefficients estimés.
\end{definition}

Le résidu des moindres carrés est la combinaison de l'échantillon qui a l'air le plus stationnaire, de sorte qu'elle paraît stationnaire plus souvent qu'une combinaison fixe~: sa statistique de Dickey--Fuller
est décalée vers la gauche, et les valeurs critiques dépendent du nombre de variables. Pour trois variables, la valeur à 5~\% de MacKinnon (2010) est $-3.74$, contre
$-2.86$ pour une seule série~; dans la simulation du chapitre portant sur 4\,000 triplets de marches aléatoires indépendantes, utiliser la valeur d'une seule série détecterait une \omterm{def:qm:multivariate-series-and-cointegration:coint}{cointégration} dans 27~\% des
cas (\cref{fig:qm:multivariate-series-and-cointegration:egdf}). Régresser le rendement à cinq ans sur les deux autres donne 0.40 fois le deux ans plus 0.62 fois
le dix ans, et une statistique résiduelle de $-10.04$~: \omterm{def:qm:multivariate-series-and-cointegration:coint}{cointégration}, très nettement.

\begin{omfigure}
\begin{tikzpicture}
\begin{axis}[width=10cm,height=5cm,xlabel={statistique $t$ de Dickey--Fuller},ylabel={densité sous l'hypothèse nulle},
  tick label style={font=\small},label style={font=\small},xmin=-6,xmax=2,ymin=0,ymax=0.6,grid=major,grid style={gray!25},
  legend columns=2,legend style={font=\footnotesize,at={(0.5,-0.3)},anchor=north,draw=none,fill=none,
  /tikz/every even column/.append style={column sep=8pt}},table/col sep=comma]
\addplot[omThm,thick,const plot mark mid,fill=omThm!15] table[x=tau,y=df]{figdata/methods/20-multivariate-series-and-cointegration/egdf.csv};
\addplot[omDef,very thick,const plot mark mid] table[x=tau,y=eg]{figdata/methods/20-multivariate-series-and-cointegration/egdf.csv};
\addplot[omThm,thin,dashed,sharp plot] coordinates {(-2.92,0) (-2.92,0.6)};
\addplot[omDef,thin,dashed,sharp plot] coordinates {(-3.75,0) (-3.75,0.6)};
\legend{une marche aléatoire,résidu d'Engle--Granger (3 séries)}
\end{axis}
\end{tikzpicture}

\omcaption{Distributions nulles de la statistique de Dickey--Fuller pour une marche aléatoire et pour le résidu des moindres carrés d'une marche aléatoire régressée sur deux
autres, à partir de 4\,000 échantillons simulés de 1\,000 jours~; les traits en pointillés correspondent à leurs points à 5~\%, $-2.92$ et $-3.75$. Données~: le tutoriel du chapitre,
à graine fixée.}\label{fig:qm:multivariate-series-and-cointegration:egdf}
\end{omfigure}

\section{Tests de rang}

\begin{definition}[Test de Johansen]\label{def:qm:multivariate-series-and-cointegration:johansen}
Le \emph{test de Johansen}\index{test de Johansen} (Johansen, 1988, 1991) estime le \omterm{def:qm:multivariate-series-and-cointegration:vecm}{modèle vectoriel à correction d'erreur} (VECM) par maximum de vraisemblance sous la contrainte $\operatorname{rank}\Pi = r$~:
après élimination de la dynamique de court terme de $\Delta Y_t$ et $Y_{t-1}$, les carrés des corrélations canoniques $\hat\lambda_1 \ge \dots \ge \hat\lambda_k$ entre les deux
ensembles de résidus donnent la statistique de la trace $-T\sum_{i>r}\ln(1 - \hat\lambda_i)$ pour un rang au plus égal à $r$, et la statistique de la valeur propre maximale $-T\ln(1 - \hat\lambda_{r+1})$~;
les vecteurs propres correspondants estiment $\beta$.
\end{definition}

Les distributions nulles sont des fonctionnelles du \omterm{def:qm:brownian-motion:bm}{mouvement brownien} qui dépendent du nombre de tendances communes et des termes déterministes. Avec la constante
restreinte à l'espace de \omterm{def:qm:multivariate-series-and-cointegration:coint}{cointégration} (les rendements n'ont pas de tendance), le chapitre les simule à partir de 20\,000 systèmes de marches aléatoires de 2\,000 jours~; les points à 95~\% de la
statistique de la trace sont 35.23, 20.36 et 9.23 pour trois, deux et une tendance commune.

Sur les trois rendements, 1976--2026, avec neuf différences retardées (le VAR(10) de l'AIC), les statistiques de la trace sont 59.77 pour un rang nul, 14.93 pour un rang au plus égal à un
et 2.34 pour un rang au plus égal à deux~: rang un, soit un seul \omterm{def:qm:multivariate-series-and-cointegration:coint}{vecteur de cointégration}. Normalisé sur le rendement à cinq ans, il vaut $0.41 \times$ rendement à deux ans $- 1 \times$ rendement à cinq ans $+ 0.61
\times$ rendement à dix ans~: un butterfly de courbe dont les ailes totalisent 1.02, proche d'une neutralité de niveau, mais incliné vers le dix ans, si bien qu'il neutralise aussi l'essentiel de la pente. En tant qu'AR(1),
il revient à la moyenne avec un coefficient de 0.9841, une \omterm{def:qm:stochastic-differential-equations:ou}{demi-vie} de 43 jours de bourse~; le butterfly 1-2-1 (ramené au même poids pour le rendement à cinq ans) a un coefficient de 0.9891 et une \omterm{def:qm:stochastic-differential-equations:ou}{demi-vie}
de 63 jours, et une volatilité quotidienne de 2.15 points de base contre 2.06 (la \cref{fig:qm:multivariate-series-and-cointegration:flies}).

\begin{omfigure}
\begin{tikzpicture}
\begin{groupplot}[group style={group size=1 by 2,vertical sep=0.5cm,xticklabels at=edge bottom},width=12cm,height=4.2cm,
  tick label style={font=\small},label style={font=\small},xmin=1976,xmax=2026.8,grid=major,grid style={gray!25},
  xticklabel style={/pgf/number format/set thousands separator={}},table/col sep=comma]
\nextgroupplot[ylabel={rendement (pb)},ymin=0,ymax=1700,legend columns=3,
  legend style={font=\footnotesize,at={(0.98,0.96)},anchor=north east,fill=white,draw=none}]
\addplot[omDef,thick] table[x=year,y=y2]{figdata/methods/20-multivariate-series-and-cointegration/yields.csv};
\addplot[black,thick] table[x=year,y=y5]{figdata/methods/20-multivariate-series-and-cointegration/yields.csv};
\addplot[omThm,thick] table[x=year,y=y10]{figdata/methods/20-multivariate-series-and-cointegration/yields.csv};
\legend{2 ans,5 ans,10 ans}
\nextgroupplot[xlabel={année},ylabel={butterfly (pb, centré)},ymin=-80,ymax=80,legend columns=2,
  legend style={font=\footnotesize,at={(0.98,0.96)},anchor=north east,fill=white,draw=none}]
\addplot[omThm!70,thin] table[x=year,y=fly121]{figdata/methods/20-multivariate-series-and-cointegration/yields.csv};
\addplot[omDef,thin] table[x=year,y=flyj]{figdata/methods/20-multivariate-series-and-cointegration/yields.csv};
\legend{1-2-1,Johansen}
\end{groupplot}
\end{tikzpicture}

\omcaption{En haut~: les rendements des bons du Trésor américain à 2, 5 et 10 ans, 1976--2026, chacun indiscernable d'une marche aléatoire. En bas~: deux butterflies, en position longue sur le 5 ans et courte sur les
ailes, le poids du 5 ans valant un~: les poids 1-2-1 (0.5, 0.5) et les poids de Johansen (0.41, 0.61), dont les \omterm{def:qm:stochastic-differential-equations:ou}{demi-vies} sont de 63 et 43 jours de bourse. Données~:
séries FRED DGS2, DGS5 et DGS10 (H.15).}\label{fig:qm:multivariate-series-and-cointegration:flies}
\end{omfigure}

Les poids ne sont qu'une estimation, et une estimation instable. Sur des fenêtres glissantes de cinq ans, le test de la trace ne détecte aucune \omterm{def:qm:multivariate-series-and-cointegration:coint}{cointégration} à 5~\% dans 20 des 46 fenêtres, et dans les
fenêtres où il en détecte un, le poids à deux ans varie entre 0.27 et 0.63 et le poids à dix ans entre 0.40 et 0.97
(la \cref{fig:qm:multivariate-series-and-cointegration:rolling}). Cinquante ans déterminent le butterfly~; cinq ans ne le permettent pas, et un desk qui négocie le butterfly vit dans ces cinq années.

\begin{omfigure}
\begin{tikzpicture}
\begin{axis}[width=11cm,height=5cm,xlabel={dernière année de la fenêtre de cinq ans},ylabel={poids du butterfly},
  tick label style={font=\small},label style={font=\small},xmin=1980,xmax=2027,ymin=0,ymax=1.2,grid=major,grid style={gray!25},
  xticklabel style={/pgf/number format/set thousands separator={}},
  legend columns=2,legend style={font=\footnotesize,at={(0.5,-0.3)},anchor=north,draw=none,fill=none,
  /tikz/every even column/.append style={column sep=8pt}},table/col sep=comma]
\addplot[only marks,mark=*,mark size=1.8pt,omDef] table[x=end,y=w2]{figdata/methods/20-multivariate-series-and-cointegration/rolling.csv};
\addplot[only marks,mark=square*,mark size=1.8pt,omThm] table[x=end,y=w10]{figdata/methods/20-multivariate-series-and-cointegration/rolling.csv};
\addplot[omDef,dashed,thin,sharp plot] coordinates {(1980,0.411) (2027,0.411)};
\addplot[omThm,dashed,thin,sharp plot] coordinates {(1980,0.607) (2027,0.607)};
\legend{poids à 2 ans,poids à 10 ans}
\end{axis}
\end{tikzpicture}

\omcaption{Poids du butterfly de Johansen (par unité du 5 ans) estimés sur des fenêtres glissantes de cinq ans, avec les valeurs sur l'échantillon complet, 0.41 et 0.61 (en pointillés). Quatre fenêtres
dont les poids dépassent 1.5 en valeur absolue ne sont pas représentées~; dans 20 des 46 fenêtres, le test de la trace ne détecte pas de \omterm{def:qm:multivariate-series-and-cointegration:coint}{cointégration} à 5~\%. Données~: séries FRED DGS2,
DGS5 et DGS10 (H.15).}\label{fig:qm:multivariate-series-and-cointegration:rolling}
\end{omfigure}

\section{Tutoriel~: le butterfly qui revient à la moyenne}\label{tut:qm:multivariate-series-and-cointegration:fly}

\begin{tutorial}
\textbf{Objectif.} Trouver le butterfly de \omterm{def:qm:multivariate-series-and-cointegration:coint}{cointégration} des rendements à 2, 5 et 10 ans et le comparer au butterfly 1-2-1.
\textbf{État final~:} les \cref{fig:qm:multivariate-series-and-cointegration:flies,fig:qm:multivariate-series-and-cointegration:rolling}, rang un, les poids (0.41,
0.61) et les \omterm{def:qm:stochastic-differential-equations:ou}{demi-vies} de 43 et 63 jours.
\begin{tutsteps}
\item \textbf{Régression de Johansen à rang réduit}, avec la constante restreinte à l'espace de \omterm{def:qm:multivariate-series-and-cointegration:coint}{cointégration}.
\omcode{code/firm/coint/firm_coint.py}{165}{198}{Corrélations canoniques de Johansen, statistiques de la trace et de la valeur propre maximale.}\label{lst:qm:multivariate-series-and-cointegration:johansen}
\item \textbf{Engle--Granger} avec les surfaces de réponse de MacKinnon.
\omcode{code/firm/coint/firm_coint.py}{146}{162}{Test d'Engle--Granger sur les résidus.}\label{lst:qm:multivariate-series-and-cointegration:eg}
\item \textbf{Exécuter} \texttt{cointegration()}, \texttt{var\_analysis()}, \texttt{rolling\_weights()} et \texttt{eg\_vs\_df()} dans \texttt{qm\_coint.py}, puis
\texttt{fig\_coint.py}.
\end{tutsteps}
\textbf{Pour aller plus loin.} Pondérer le butterfly par la duration (DV01) plutôt que par le rendement et refaire le test sur le P\&L~; estimer le coefficient $\alpha$ du VECM et voir quel
rendement assure la correction.
\end{tutorial}

\section{Construction~: le module de cointégration}\label{bld:qm:multivariate-series-and-cointegration:coint}

\begin{build}
\bfield{Rôle} Les paniers de valeur relative de la firme miniature (butterflies, spreads, indice contre ses composantes) reçoivent ici leurs poids et leurs tests.
\bfield{Interface} \texttt{var\_fit(Y, p)}, \texttt{var\_stable(A)}, \texttt{irf(A, Sigma, horizon)}, \texttt{fevd(A, Sigma, horizon)}~;
\texttt{granger\_test(Y, p, cause, effect)}, \texttt{f\_sf}~; \texttt{engle\_granger(y, X, lags)}~; \texttt{johansen(Y, lags)}~; \texttt{johansen\_crit\_sim(m, T, reps, seed)}.
\bfield{Règles} Les valeurs critiques proviennent des surfaces de réponse ou de la simulation documentée de ce module, jamais de la table de Dickey--Fuller~; les poids sont reportés
avec le rang et l'échantillon~; l'ordre des variables est précisé pour chaque réponse orthogonalisée.
\bfield{Tests d'acceptation} \texttt{code/firm/coint/tests/}~: un VAR(1) retrouvé, stabilité, les premières \omterm{def:qm:multivariate-series-and-cointegration:irf}{fonctions de réponse impulsionnelle}, décompositions dont la somme vaut un~; tests de Granger avec
puissance et taille correcte~; la queue de $F$ face à la limite en $\chi^2$~; Engle--Granger sur une paire cointégrée et sur des marches indépendantes~; Johansen trouve un rang un et sa
direction~; les valeurs critiques stockées sont re-simulées.
\bfield{Pour aller plus loin} Les $\alpha$ et $\Gamma$ du VECM avec leurs \omterm{def:qm:estimation:estimator}{erreurs types}~; tests d'exogénéité faible~; le butterfly négocié avec un poids filtré par Kalman (chapitre~19).
\end{build}

\begin{omsources}
\item C.~A.~Sims, ``Macroeconomics and reality'', \emph{Econometrica} 48, 1980.
\item C.~W.~J.~Granger, ``Investigating causal relations by econometric models and cross-spectral methods'', \emph{Econometrica} 37, 1969.
\item R.~F.~Engle and C.~W.~J.~Granger, ``Co-integration and error correction: representation, estimation, and testing'', \emph{Econometrica} 55, 1987.
\item S.~Johansen, ``Statistical analysis of cointegration vectors'', \emph{Journal of Economic Dynamics and Control} 12, 1988; ``Estimation and hypothesis
testing of cointegration vectors in Gaussian vector autoregressive models'', \emph{Econometrica} 59, 1991.
\item J.~G.~MacKinnon, ``Critical values for cointegration tests'', Queen's Economics Department Working Paper 1227, 2010.
\item Board of Governors of the Federal Reserve System, H.15 Selected Interest Rates, via FRED (séries DGS2, DGS5 et DGS10), consulté le 24 septembre 2026.
\end{omsources}

\section{Exercices}

\begin{exercise}[$\star$]\label{exo:qm:multivariate-series-and-cointegration:1}
Le VAR(1) de matrice $A = \begin{psmallmatrix}0.5 & 0.4\\ 0.4 & 0.5\end{psmallmatrix}$ est-il stationnaire~? Et celui de matrice $\begin{psmallmatrix}0.6 & 0.4\\ 0.4 & 0.6\end{psmallmatrix}$~?
\end{exercise}

\begin{exercise}[$\star$]\label{exo:qm:multivariate-series-and-cointegration:2}
Deux rendements $x$ et $y$ sont des marches aléatoires, et $y - 0.9x$ est stationnaire. Quels sont $k$, $r$ et le nombre de tendances communes~?
\end{exercise}

\begin{exercise}[$\star$]\label{exo:qm:multivariate-series-and-cointegration:3}
Calculer la valeur critique à 5~\% d'Engle--Granger de MacKinnon pour trois variables et $T = 500$, à partir de $(-3.74066, -8.5631, -10.852, 27.982)$.
\end{exercise}

\begin{exercise}[$\star\star$]\label{exo:qm:multivariate-series-and-cointegration:4}
Écrire le modèle VECM de $y_t = 0.9x_{t-1} + 0.5(y_{t-1} - 0.9x_{t-1}) + e_t$ avec $x$ une marche aléatoire, et identifier $\alpha$ et $\beta$.
\end{exercise}

\begin{exercise}[$\star\star$]\label{exo:qm:multivariate-series-and-cointegration:5}
Deux résidus ont une corrélation de 0.9 et des variances unitaires. La première variable étant ordonnée en premier, quelle fraction de la variance de la prévision à un pas de la seconde variable est
attribuée au premier choc~?
\end{exercise}

\begin{exercise}[$\star\star$]\label{exo:qm:multivariate-series-and-cointegration:6}
Les valeurs propres de Johansen d'un système à trois variables avec $T = 1\,000$ sont 0.05, 0.01 et 0.002. Calculer les statistiques de la trace et le rang à 5~\%.
\end{exercise}

\begin{exercise}[$\star\star\star$]\label{exo:qm:multivariate-series-and-cointegration:7}
\emph{Programmation.} Estimer les poids de Johansen du fly sur des fenêtres glissantes de cinq ans de rendements, puis compter les fenêtres dans lesquelles le test de la trace ne détecte aucune \omterm{def:qm:multivariate-series-and-cointegration:coint}{cointégration}.
\end{exercise}

\begin{exercise}[$\star\star\star$]\label{exo:qm:multivariate-series-and-cointegration:8}
\emph{Repérer l'erreur.} «~Nous avons régressé le rendement à 5 ans sur ceux à 2 ans et à 10 ans, obtenu un résidu dont la statistique de Dickey--Fuller vaut $-3.1$, en dessous de la valeur critique à 5~\%
de $-2.86$, et nous avons commencé à trader le résidu comme un fly stationnaire~».
\end{exercise}

\section{Problème~: le fly qui revient à la moyenne}

\begin{problem}[{Problème du week-end --- le fly 2s5s10s cointégré}]\label{pb:qm:multivariate-series-and-cointegration:1}
Les données sont les rendements quotidiens à 2, 5 et 10 ans des bons du Trésor à maturité constante (FRED DGS2, DGS5, DGS10), du 1\up{er} juin 1976 au 22 septembre 2026, soit 12\,574 jours, exprimés en points de base.

\medskip\noindent\textbf{Partie I --- Chaque série.}
\begin{enumerate}
\item Que disent les \omterm{def:qm:linear-time-series:df}{tests de Dickey--Fuller} sur chaque rendement~?
\item Quel ordre de VAR le critère d'AIC retient-il pour les variations quotidiennes, et le VAR est-il stable~?
\item Quels sont les écarts-types et les corrélations des résidus~?
\item Que produit un choc d'un écart-type sur le rendement à 2 ans pendant dix jours~?
\item Quels rendements causent au sens de Granger quels autres, et que signifie cette relation~?
\end{enumerate}

\medskip\noindent\textbf{Partie II --- \omterm{def:qm:multivariate-series-and-cointegration:coint}{Cointégration}.}
\begin{enumerate}[resume]
\item Quelles sont les statistiques de la trace de Johansen et le rang à 5~\%~?
\item Quels sont les poids du fly par unité du rendement à 5 ans, et que neutralisent-ils~?
\item Que donne Engle--Granger, et par rapport à quelle valeur critique~?
\item Quelles sont les \omterm{def:qm:stochastic-differential-equations:ou}{demi-vies} du fly de Johansen et du fly un-deux-un~?
\item Quelles sont les volatilités quotidiennes des flies~?
\end{enumerate}

\medskip\noindent\textbf{Partie III --- Robustesse.}
\begin{enumerate}[resume]
\item À quelle fréquence une fenêtre de cinq ans détecte-t-elle une \omterm{def:qm:multivariate-series-and-cointegration:coint}{cointégration}~?
\item De combien les poids varient-ils d'une fenêtre à l'autre~?
\item Pourquoi la valeur critique d'Engle--Granger est-elle plus basse que celle de Dickey--Fuller~?
\item Comment corriger la \omterm{def:qm:stochastic-differential-equations:ou}{demi-vie} du biais de petit échantillon du chapitre~17~?
\item Que risquerait un desk en négociant les poids estimés sur l'échantillon complet~?
\end{enumerate}

\medskip\noindent\textbf{Partie IV --- Jugement.}
\begin{enumerate}[resume]
\item Quel fly négocier, et avec quels poids~?
\item Comment le dimensionner pour une \omterm{def:qm:stochastic-differential-equations:ou}{demi-vie} de 43 jours~?
\item Pourquoi faut-il réestimer les poids, et à quelle fréquence~?
\item Énoncer le \emph{résultat nommé}~: les poids et la \omterm{def:qm:stochastic-differential-equations:ou}{demi-vie} du fly de Johansen comparés à ceux du fly un-deux-un.
\item En une phrase~: que promet la \omterm{def:qm:multivariate-series-and-cointegration:coint}{cointégration} à un trader en valeur relative, et que ne promet-elle pas~?
\end{enumerate}
\end{problem}

\section{Questions d'entretien}

\begin{interviewq}[$\star$ \iqroles{chercheur, trader}]\label{iq:qm:multivariate-series-and-cointegration:1}
Quelle est la différence entre corrélation et \omterm{def:qm:multivariate-series-and-cointegration:coint}{cointégration}~?
\end{interviewq}

\begin{interviewq}[$\star\star$ \iqroles{chercheur, trader}]\label{iq:qm:multivariate-series-and-cointegration:2}
Comment choisiriez-vous les poids d'un fly de courbe pour négocier le \omterm{def:qm:stochastic-differential-equations:ou}{retour à la moyenne}~?
\end{interviewq}

\begin{interviewq}[$\star\star$ \iqroles{chercheur}]\label{iq:qm:multivariate-series-and-cointegration:3}
Pourquoi ne peut-on pas utiliser les valeurs critiques de Dickey--Fuller sur le résidu d'une régression de \omterm{def:qm:multivariate-series-and-cointegration:coint}{cointégration}~?
\end{interviewq}

\begin{interviewq}[$\star\star$ \iqroles{chercheur, ingénieur ML}]\label{iq:qm:multivariate-series-and-cointegration:4}
Que teste la causalité au sens de Granger, et que ne teste-t-elle pas~?
\end{interviewq}

\begin{interviewq}[$\star\star$ \iqroles{chercheur}]\label{iq:qm:multivariate-series-and-cointegration:5}
Pourquoi l'ordre des variables compte-t-il dans les réponses impulsionnelles d'un VAR~?
\end{interviewq}

\begin{interviewq}[$\star\star\star$ \iqroles{chercheur}]\label{iq:qm:multivariate-series-and-cointegration:6}
Expliquer le \omterm{def:qm:multivariate-series-and-cointegration:johansen}{test de Johansen} en termes de corrélations canoniques.
\end{interviewq}
